\documentclass[11pt]{article}

\usepackage[margin=1in]{geometry}
\usepackage{amsmath,amssymb}
\usepackage{hyperref}
\usepackage{enumitem}

\hypersetup{colorlinks=true,linkcolor=blue,urlcolor=blue,citecolor=blue}
\emergencystretch=2em

\title{The $\mathrm{RCD}(K,\infty)$ Log-Sobolev Rigidity Problem\\
Was Answered by Han}
\author{Math discovery packet}
\date{August 11, 2026}

\begin{document}
\maketitle

\section*{Status and source problem}

This is a lightweight \emph{literature-already-answered} status note, not an
original proof.  Ohta and Takatsu, arXiv:1904.09400, Section 4 (``Further
problems''), item (B), page 10 of the locally compiled source PDF, write:
\begin{quote}
Cheng--Zhou's characterization of the sharp spectral gap was extended to
$\mathrm{RCD}(K,\infty)$-spaces.  Hence it is natural to expect the analogue
to Theorem 1.1 for $\mathrm{RCD}(K,\infty)$-spaces.  They suggest either
extending needle decomposition to this setting or directly showing that
$\log\rho$ attains the sharp spectral gap.
\end{quote}
Thus the requested analogue says that a nonconstant equality case in the
sharp logarithmic Sobolev inequality, under $K>0$, forces a one-dimensional
Gaussian factor and that the extremal density is a translation along that
factor.

\section*{The later answer}

Bang-Xian Han, \emph{Rigidity of some functional inequalities on RCD
spaces}, arXiv:2001.07930, Theorem 1.7 (page 8 of the locally compiled
supporting PDF), proves the following.  If an $\mathrm{RCD}(K,\infty)$
probability space with $K>0$ admits a nonconstant equality case for an
admissible $\Phi$-entropy inequality, then it splits isometrically as
\[
 \left(\mathbb R,|\cdot|,
 \sqrt{\frac{K}{2\pi}}e^{-Kr^2/2}\,dr\right)
 \times (Y,d_Y,m_Y).
\]
For $\Phi(x)=x\log x$, the logarithmic-Sobolev extremal is
\[
 f(r,y)=\exp\!\left(ar-\frac{a^2}{2K}\right),
 \qquad a\in\mathbb R,
\]
which is precisely the Radon--Nikodym density of a translation along the
Gaussian coordinate.  This gives both conclusions of the desired analogue:
Gaussian splitting and classification of equality functions.

\section*{Why the provenance is explicit}

Han's Remark 1.9 (page 9) cites Ohta--Takatsu and states that their Section 4
left rigidity of the log-Sobolev inequality on
$\mathrm{RCD}(K,\infty)$ spaces open because dimension-free needle
decomposition was unavailable.  The remark then says explicitly that Han's
result extends Ohta--Takatsu to this setting.  The supporting author therefore
knew that the theorem answered the same problem.

\section*{Scope}

The status is \texttt{literature\_already\_answered (Further problem (B),
affirmative)}.  Items (A), (C), (D), and (E) in the source's Section 4 are not
assessed here.

\section*{References}

\begin{itemize}[leftmargin=*]
\item S.-i. Ohta and A. Takatsu, \emph{Equality in the logarithmic Sobolev
inequality}, arXiv:1904.09400, Section 4(B).
\item B.-X. Han, \emph{Rigidity of some functional inequalities on RCD
spaces}, arXiv:2001.07930, Theorem 1.7 and Remark 1.9.
\end{itemize}

\end{document}
